\begin{thebibliography}{57}
\bibitem{bullmore2009}
E. Bullmore and O. Sporns, ``Complex brain networks: Graph theoretical analysis of structural and functional systems,'' \emph{Nat. Rev. Neurosci.}, vol. 10, no. 3, pp. 186--198, Mar. 2009, doi: 10.1038/nrn2575.

\bibitem{smith2011}
S. M. Smith \emph{et al.}, ``Network modelling methods for FMRI,'' \emph{NeuroImage}, vol. 54, no. 2, pp. 875--891, Jan. 2011, doi: 10.1016/j.neuroimage.2010.08.063.

\bibitem{dimartino2014}
A. Di Martino \emph{et al.}, ``The autism brain imaging data exchange: Towards a large-scale evaluation of the intrinsic brain architecture in autism,'' \emph{Mol. Psychiatry}, vol. 19, no. 6, pp. 659--667, Jun. 2014, doi: 10.1038/mp.2013.78.

\bibitem{parisot2018}
S. Parisot \emph{et al.}, ``Disease prediction using graph convolutional networks: Application to autism spectrum disorder and Alzheimer's disease,'' \emph{Med. Image Anal.}, vol. 48, pp. 117--130, Aug. 2018, doi: 10.1016/j.media.2018.06.001.

\bibitem{li2021braingnn}
X. Li \emph{et al.}, ``BrainGNN: Interpretable brain graph neural network for fMRI analysis,'' \emph{Med. Image Anal.}, vol. 74, Art. no. 102233, Dec. 2021, doi: 10.1016/j.media.2021.102233.

\bibitem{kan2022}
X. Kan, W. Dai, H. Cui, Z. Zhang, Y. Guo, and C. Yang, ``Brain network transformer,'' in \emph{Adv. Neural Inf. Process. Syst. (NeurIPS)}, vol. 35, 2022.

\bibitem{cui2023}
H. Cui \emph{et al.}, ``BrainGB: A benchmark for brain network analysis with graph neural networks,'' \emph{IEEE Trans. Med. Imag.}, vol. 42, no. 2, pp. 493--506, Feb. 2023, doi: 10.1109/TMI.2022.3218745.

\bibitem{abraham2017}
A. Abraham \emph{et al.}, ``Deriving reproducible biomarkers from multi-site resting-state data: An autism-based example,'' \emph{NeuroImage}, vol. 147, pp. 736--745, Feb. 2017, doi: 10.1016/j.neuroimage.2016.10.045.

\bibitem{heinsfeld2018}
A. S. Heinsfeld, A. R. Franco, R. C. Craddock, A. Buchweitz, and F. Meneguzzi, ``Identification of autism spectrum disorder using deep learning and the ABIDE dataset,'' \emph{NeuroImage Clin.}, vol. 17, pp. 16--23, 2018, doi: 10.1016/j.nicl.2017.08.017.

\bibitem{yu2018harm}
M. Yu \emph{et al.}, ``Statistical harmonization corrects site effects in functional connectivity measurements from multi-site fMRI data,'' \emph{Hum. Brain Mapp.}, vol. 39, no. 11, pp. 4213--4227, Nov. 2018, doi: 10.1002/hbm.24241.

\bibitem{hutchison2013}
R. M. Hutchison \emph{et al.}, ``Dynamic functional connectivity: Promise, issues, and interpretations,'' \emph{NeuroImage}, vol. 80, pp. 360--378, Oct. 2013, doi: 10.1016/j.neuroimage.2013.05.079.

\bibitem{allen2014}
E. A. Allen, E. Damaraju, S. M. Plis, E. B. Erhardt, T. Eichele, and V. D. Calhoun, ``Tracking whole-brain connectivity dynamics in the resting state,'' \emph{Cereb. Cortex}, vol. 24, no. 3, pp. 663--676, Mar. 2014, doi: 10.1093/cercor/bhs352.

\bibitem{preti2017}
M. G. Preti, T. A. Bolton, and D. Van De Ville, ``The dynamic functional connectome: State-of-the-art and perspectives,'' \emph{NeuroImage}, vol. 160, pp. 41--54, Oct. 2017, doi: 10.1016/j.neuroimage.2016.12.061.

\bibitem{kim2021}
B.-H. Kim, J. C. Ye, and J.-J. Kim, ``Learning dynamic graph representation of brain connectome with spatio-temporal attention,'' in \emph{Adv. Neural Inf. Process. Syst. (NeurIPS)}, vol. 34, 2021, arXiv:2105.13495.

\bibitem{gadgil2020}
S. Gadgil, Q. Zhao, A. Pfefferbaum, E. V. Sullivan, E. Adeli, and K. M. Pohl, ``Spatio-temporal graph convolution for resting-state fMRI analysis,'' in \emph{Med. Image Comput. Comput.-Assist. Intervent. (MICCAI)}, 2020, pp. 528--538, doi: 10.1007/978-3-030-59728-3\_52.

\bibitem{marquand2016}
A. F. Marquand, I. Rezek, J. Buitelaar, and C. F. Beckmann, ``Understanding heterogeneity in clinical cohorts using normative models: Beyond case-control studies,'' \emph{Biol. Psychiatry}, vol. 80, no. 7, pp. 552--561, Oct. 2016, doi: 10.1016/j.biopsych.2015.12.023.

\bibitem{ma2023survey}
X. Ma \emph{et al.}, ``A comprehensive survey on graph anomaly detection with deep learning,'' \emph{IEEE Trans. Knowl. Data Eng.}, vol. 35, no. 12, pp. 12012--12038, Dec. 2023, doi: 10.1109/TKDE.2021.3118815.

\bibitem{gotham2009}
K. Gotham, A. Pickles, and C. Lord, ``Standardizing ADOS scores for a measure of severity in autism spectrum disorders,'' \emph{J. Autism Dev. Disord.}, vol. 39, no. 5, pp. 693--705, May 2009, doi: 10.1007/s10803-008-0674-3.

\bibitem{xu2019gin}
K. Xu, W. Hu, J. Leskovec, and S. Jegelka, ``How powerful are graph neural networks?'' in \emph{Proc. Int. Conf. Learn. Represent. (ICLR)}, 2019, arXiv:1810.00826.

\bibitem{cho2014}
K. Cho \emph{et al.}, ``Learning phrase representations using RNN encoder-decoder for statistical machine translation,'' in \emph{Proc. Conf. Empirical Methods Natural Lang. Process. (EMNLP)}, Doha, Qatar, Oct. 2014, pp. 1724--1734, doi: 10.3115/v1/D14-1179.

\bibitem{glasser2016}
M. F. Glasser \emph{et al.}, ``A multi-modal parcellation of human cerebral cortex,'' \emph{Nature}, vol. 536, no. 7615, pp. 171--178, Aug. 2016, doi: 10.1038/nature18933.

\bibitem{ktena2017}
S. I. Ktena \emph{et al.}, ``Distance metric learning using graph convolutional networks: Application to functional brain networks,'' in \emph{Med. Image Comput. Comput.-Assist. Intervent. (MICCAI)}, 2017, pp. 469--477, doi: 10.1007/978-3-319-66182-7\_54.

\bibitem{kawahara2017}
J. Kawahara \emph{et al.}, ``BrainNetCNN: Convolutional neural networks for brain networks; towards predicting neurodevelopment,'' \emph{NeuroImage}, vol. 146, pp. 1038--1049, Feb. 2017, doi: 10.1016/j.neuroimage.2016.09.046.

\bibitem{kipf2017}
T. N. Kipf and M. Welling, ``Semi-supervised classification with graph convolutional networks,'' in \emph{Proc. Int. Conf. Learn. Represent. (ICLR)}, 2017, arXiv:1609.02907.

\bibitem{lanciano2020}
T. Lanciano, F. Bonchi, and A. Gionis, ``Explainable classification of brain networks via contrast subgraphs,'' in \emph{Proc. ACM SIGKDD Int. Conf. Knowl. Discovery Data Mining (KDD)}, 2020, pp. 3308--3318, doi: 10.1145/3394486.3403383.

\bibitem{rossi2020tgn}
E. Rossi, B. Chamberlain, F. Frasca, D. Eynard, F. Monti, and M. Bronstein, ``Temporal graph networks for deep learning on dynamic graphs,'' arXiv:2006.10637, Jun. 2020.

\bibitem{pareja2020}
A. Pareja \emph{et al.}, ``EvolveGCN: Evolving graph convolutional networks for dynamic graphs,'' in \emph{Proc. AAAI Conf. Artif. Intell.}, vol. 34, no. 4, 2020, pp. 5363--5370, doi: 10.1609/aaai.v34i04.5984.

\bibitem{aggarwal2019}
P. Aggarwal and A. Gupta, ``Multivariate graph learning for detecting aberrant connectivity of dynamic brain networks in autism,'' \emph{Med. Image Anal.}, vol. 56, pp. 11--25, Aug. 2019, doi: 10.1016/j.media.2019.05.007.

\bibitem{pang2021review}
G. Pang, C. Shen, L. Cao, and A. van den Hengel, ``Deep learning for anomaly detection: A review,'' \emph{ACM Comput. Surv.}, vol. 54, no. 2, 2022, doi: 10.1145/3439950.

\bibitem{ding2019dominant}
K. Ding, J. Li, R. Bhanushali, and H. Liu, ``Deep anomaly detection on attributed networks,'' in \emph{Proc. SIAM Int. Conf. Data Mining (SDM)}, 2019, pp. 594--602, doi: 10.1137/1.9781611975673.67.

\bibitem{ruff2018}
L. Ruff \emph{et al.}, ``Deep one-class classification,'' in \emph{Proc. Int. Conf. Mach. Learn. (ICML)}, PMLR vol. 80, 2018, pp. 4393--4402.

\bibitem{ma2022glocal}
R. Ma, G. Pang, L. Chen, and A. van den Hengel, ``Deep graph-level anomaly detection by glocal knowledge distillation,'' in \emph{Proc. ACM Int. Conf. Web Search Data Mining (WSDM)}, 2022, pp. 704--714, doi: 10.1145/3488560.3498473.

\bibitem{yu2018netwalk}
W. Yu, W. Cheng, C. C. Aggarwal, K. Zhang, H. Chen, and W. Wang, ``NetWalk: A flexible deep embedding approach for anomaly detection in dynamic networks,'' in \emph{Proc. ACM SIGKDD Int. Conf. Knowl. Discovery Data Mining (KDD)}, 2018, pp. 2672--2681, doi: 10.1145/3219819.3220024.

\bibitem{zheng2019addgraph}
L. Zheng, Z. Li, J. Li, Z. Li, and J. Gao, ``AddGraph: Anomaly detection in dynamic graph using attention-based temporal GCN,'' in \emph{Proc. Int. Joint Conf. Artif. Intell. (IJCAI)}, 2019, pp. 4419--4425, doi: 10.24963/ijcai.2019/614.

\bibitem{behrouz2023}
A. Behrouz and M. Seltzer, ``Anomaly detection in human brain via inductive learning on temporal multiplex networks,'' in \emph{Proc. Mach. Learn. Healthcare Conf.}, PMLR vol. 219, 2023, pp. 50--75.

\bibitem{sadeghian2024}
S. Sadeghian, X. Li, and M. Seltzer, ``HyperBrain: Anomaly detection for temporal hypergraph brain networks,'' arXiv:2410.02087, Oct. 2024.

\bibitem{chatterjee2021}
T. Chatterjee, R. Albert, S. Thapliyal, N. Azarhooshang, and B. DasGupta, ``Detecting network anomalies using Forman-Ricci curvature and a case study for human brain networks,'' \emph{Sci. Rep.}, vol. 11, no. 1, Art. no. 8121, Apr. 2021, doi: 10.1038/s41598-021-87587-z.

\bibitem{wang2021fewshot}
Y. Wang, Q. Yao, J. T. Kwok, and L. M. Ni, ``Generalizing from a few examples: A survey on few-shot learning,'' \emph{ACM Comput. Surv.}, vol. 53, no. 3, Art. no. 63, 2020, doi: 10.1145/3386252.

\bibitem{finn2017}
C. Finn, P. Abbeel, and S. Levine, ``Model-agnostic meta-learning for fast adaptation of deep networks,'' in \emph{Proc. Int. Conf. Mach. Learn. (ICML)}, PMLR vol. 70, 2017, pp. 1126--1135.

\bibitem{snell2017}
J. Snell, K. Swersky, and R. Zemel, ``Prototypical networks for few-shot learning,'' in \emph{Adv. Neural Inf. Process. Syst. (NIPS)}, vol. 30, 2017, arXiv:1703.05175.

\bibitem{ding2021fewshot}
K. Ding, Q. Zhou, H. Tong, and H. Liu, ``Few-shot network anomaly detection via cross-network meta-learning,'' in \emph{Proc. Web Conf. (WWW)}, 2021, pp. 2448--2456, doi: 10.1145/3442381.3449922.

\bibitem{pang2019devnet}
G. Pang, C. Shen, and A. van den Hengel, ``Deep anomaly detection with deviation networks,'' in \emph{Proc. ACM SIGKDD Int. Conf. Knowl. Discovery Data Mining (KDD)}, 2019, pp. 353--362, doi: 10.1145/3292500.3330871.

\bibitem{yang2022meta}
Y. Yang \emph{et al.}, ``Data-efficient brain connectome analysis via multi-task meta-learning,'' in \emph{Proc. ACM SIGKDD Int. Conf. Knowl. Discovery Data Mining (KDD)}, 2022, pp. 4743--4751, doi: 10.1145/3534678.3542680.

\bibitem{wei2025}
X. Wei, K. Zhao, Y. Jiao, L. He, and Y. Zhang, ``A brain graph foundation model: Pre-training and prompt-tuning across broad atlases and disorders,'' arXiv:2506.02044, 2025.

\bibitem{dhinagar2023}
N. J. Dhinagar, V. Santhalingam, K. E. Lawrence, E. Laltoo, and P. M. Thompson, ``Few-shot classification of autism spectrum disorder using site-agnostic meta-learning and brain MRI,'' in \emph{Proc. Annu. Int. Conf. IEEE Eng. Med. Biol. Soc. (EMBC)}, 2023, doi: 10.1109/EMBC40787.2023.10340852.

\bibitem{zhao2021smote}
T. Zhao, X. Zhang, and S. Wang, ``GraphSMOTE: Imbalanced node classification on graphs with graph neural networks,'' in \emph{Proc. ACM Int. Conf. Web Search Data Mining (WSDM)}, 2021, pp. 833--841, doi: 10.1145/3437963.3441720.

\bibitem{chen2020gaan}
Z. Chen, B. Liu, M. Wang, P. Dai, J. Lv, and L. Bo, ``Generative adversarial attributed network anomaly detection,'' in \emph{Proc. ACM Int. Conf. Inf. Knowl. Manage. (CIKM)}, 2020, pp. 1989--1992, doi: 10.1145/3340531.3412070.

\bibitem{zugner2018}
D. Z{\"u}gner, A. Akbarnejad, and S. G{\"u}nnemann, ``Adversarial attacks on neural networks for graph data,'' in \emph{Proc. ACM SIGKDD Int. Conf. Knowl. Discovery Data Mining (KDD)}, 2018, pp. 2847--2856, doi: 10.1145/3219819.3220078.

\bibitem{dai2018}
H. Dai \emph{et al.}, ``Adversarial attack on graph structured data,'' in \emph{Proc. Int. Conf. Mach. Learn. (ICML)}, PMLR vol. 80, 2018, pp. 1115--1124.

\bibitem{xu2019topology}
K. Xu \emph{et al.}, ``Topology attack and defense for graph neural networks: An optimization perspective,'' in \emph{Proc. Int. Joint Conf. Artif. Intell. (IJCAI)}, 2019, pp. 3961--3967, doi: 10.24963/ijcai.2019/550.

\bibitem{madry2018}
A. Madry, A. Makelov, L. Schmidt, D. Tsipras, and A. Vladu, ``Towards deep learning models resistant to adversarial attacks,'' in \emph{Proc. Int. Conf. Learn. Represent. (ICLR)}, 2018, arXiv:1706.06083.

\bibitem{brahim2020}
A. Brahim and N. Farrugia, ``Graph Fourier transform of fMRI temporal signals based on an averaged structural connectome for the classification of neuroimaging,'' \emph{Artif. Intell. Med.}, vol. 106, Art. no. 101870, 2020, doi: 10.1016/j.artmed.2020.101870.

\bibitem{abraham2014}
A. Abraham \emph{et al.}, ``Machine learning for neuroimaging with scikit-learn,'' \emph{Front. Neuroinform.}, vol. 8, Art. no. 14, 2014, doi: 10.3389/fninf.2014.00014.

\bibitem{bardes2022}
A. Bardes, J. Ponce, and Y. LeCun, ``VICReg: Variance-invariance-covariance regularization for self-supervised learning,'' in \emph{Proc. Int. Conf. Learn. Represent. (ICLR)}, 2022, arXiv:2105.04906.

\bibitem{paszke2019}
A. Paszke \emph{et al.}, ``PyTorch: An imperative style, high-performance deep learning library,'' in \emph{Adv. Neural Inf. Process. Syst. (NeurIPS)}, vol. 32, 2019, arXiv:1912.01703.

\bibitem{fey2019}
M. Fey and J. E. Lenssen, ``Fast graph representation learning with PyTorch Geometric,'' in \emph{ICLR Workshop Representation Learn. Graphs Manifolds}, 2019, arXiv:1903.02428.

\bibitem{kingma2015}
D. P. Kingma and J. Ba, ``Adam: A method for stochastic optimization,'' in \emph{Proc. Int. Conf. Learn. Represent. (ICLR)}, 2015, arXiv:1412.6980.

\end{thebibliography}
